\section{Discrete propagation inequalities}\label{sec:discrete-propagation}

We now derive the one-step inequality used to propagate the contact from
Proposition~\ref{prop:envelope-contact}. The argument will also be applied
after normalizing at fitted points that need not be contacts, so the
estimates must hold without assuming contact at the center.
Stopped stochastic processes and finite-difference comparisons have also
been used to study the infinity Laplacian
\cite{PeresSchrammSheffieldWilson2009,ArmstrongSmart2010}.
The walk constructed here uses forward and backward quadratic optimizers
for the limiting equation; its drift controls velocity discrepancies
that persist under fitted changes of frame.

Throughout this section, let $v\in\Aclass_{m,C_0}$ be centered at $(0,1)$,
and fix $0\le\kappa<\infty$ for which \eqref{eq:jensen-bound} holds with
$f=v$ and $I=\R$.  For $t>0$, define
\[
  W(x,t)=\Cenv_t v(x,t)-\frac{\kappa^2t}{2}.
\]
The general bound \eqref{eq:barycentric-sandwich}, with its parameter equal
to $t$, gives
\[
  v(x,t)-\frac{\kappa^2t}{2}\le W(x,t)\le v(x,t).
\]
Lemma~\ref{lem:barycentric-envelope} gives continuity of $W$ and concavity of
$W(\cdot,t)-|\cdot|^2/(2t)$.  The sandwich and
\eqref{eq:uniform-subquadratic-growth} give subquadratic spatial growth for
$W$ on compact positive-time intervals.  Since $v$ is a subsolution with
that growth, Corollary~\ref{cor:barycentric-subsolution}, applied with
$c=\kappa^2/2$, makes $W$ a subsolution of \eqref{eq:limit-pde}.

For $\ell,\epsilon>0$, set
\begin{equation}\label{eq:dp-auxiliary}
  f=W,\qquad f^\ell=\Penv_\ell f,\qquad
  g_\epsilon=\Menv_\epsilon v,\qquad
  J_\epsilon=\frac{f^\epsilon-g_\epsilon}{2\epsilon}.
\end{equation}
These functions are defined at every positive time.  The identity
\begin{equation}\label{eq:dp-contact-gap}
\begin{split}
  2\epsilon J_\epsilon(z)
  &=W(z)-v(z)+\bigl(\Penv_\epsilon W(z)-W(z)\bigr)\\
  &\hspace{4em}+\bigl(v(z)-\Menv_\epsilon v(z)\bigr)
\end{split}
\end{equation}
and the zero-length continuity in Lemma~\ref{lem:quadratic-envelopes} show
that $2\epsilon J_\epsilon\to W-v$ uniformly on every compact subset of
$\R^m\times(0,\infty)$.  Thus, on a compact region separated from contact,
$J_\epsilon$ tends uniformly to $-\infty$.  This fact will turn lower
expectation bounds for $J_\epsilon$ into contact at limiting marks.
The sandwich also gives $\Penv_\epsilon W\le\Penv_\epsilon v$, and hence
\begin{equation}\label{eq:dp-J-bounds}
  J_\epsilon(z)\le
  \frac12\bigl(\Senv^+_\epsilon v(z)+\Senv^-_\epsilon v(z)\bigr).
\end{equation}

For the particular pair supplied by Proposition~\ref{prop:envelope-contact},
write $z_*=(0,1)$ and $H_*=H_v(z_*)\in\R$.  Equality of $W$ and $v$ there,
the contact slope conclusion, and the common slope limit for $v$ give
\begin{equation}\label{eq:dp-J-contact}
  J_\epsilon(z_*)
  =\frac12\bigl(\Senv^+_\epsilon W(z_*)+
                    \Senv^-_\epsilon v(z_*)\bigr)
  \longrightarrow H_*\qquad(\epsilon\downarrow0).
\end{equation}
This finite limiting value will start the walks in
Section~\ref{sec:contact-localization}.

\subsection{Spatial regularity and the one-step drift}

\begin{lemma}\label{lem:dp-spatial-regularity}
For every $t,\ell>0$, the optimization defining $f^\ell(x,t)$ has a unique
maximizer $y=y(x,t,\ell)$.  The function $f^\ell(\cdot,t)$ is differentiable,
and
\begin{equation}\label{eq:dp-gradient}
  \nabla_x f^\ell(x,t)=\frac{y-x}{\ell}.
\end{equation}
For every $h\in\R^m$ it satisfies the two supporting inequalities
\begin{equation}\label{eq:dp-supports}
\begin{split}
  f^\ell(x,t)+\nabla_x f^\ell(x,t)\cdot h-\frac{|h|^2}{2\ell}
  &\le f^\ell(x+h,t)\\
  &\le f^\ell(x,t)+\nabla_x f^\ell(x,t)\cdot h+\frac{|h|^2}{2t}.
\end{split}
\end{equation}
In particular, its spatial Hessian is bounded below by $-\ell^{-1}I$ in the
semiconvexity sense and above by $t^{-1}I$ in the semiconcavity sense.
\end{lemma}

\begin{proof}
Put $s=t+\ell$.  The function $f(\cdot,s)-|\cdot|^2/(2s)$ is concave,
because the same is true of $\Cenv_s v(\cdot,s)-|\cdot|^2/(2s)$.  Thus
\[
  y\longmapsto f(y,s)-\frac{|y-x|^2}{2\ell}
\]
is strictly concave: after subtracting the concave function
$f(y,s)-|y|^2/(2s)$, its quadratic part has Hessian
$(s^{-1}-\ell^{-1})I=-t(\ell s)^{-1}I$.  It has a maximizer by the
subquadratic growth and the quadratic cost, and strict concavity makes the
maximizer unique.  Keeping this maximizer fixed when $x$ is replaced by
$x+h$ gives
\begin{equation}\label{eq:dp-lower-support}
  f^\ell(x+h,t)\ge f^\ell(x,t)+\frac{y-x}{\ell}\cdot h
                                      -\frac{|h|^2}{2\ell}.
\end{equation}

For the upper bound, the identity
\[
  \frac{|y|^2}{2s}-\frac{|y-x|^2}{2\ell}-\frac{|x|^2}{2t}
  =-\frac{t}{2\ell s}\left|y-\frac{s}{t}x\right|^2
\]
shows that
\[
  f^\ell(x,t)-\frac{|x|^2}{2t}
  =\sup_y\left\{f(y,s)-\frac{|y|^2}{2s}
                  -\frac{t}{2\ell s}\left|y-\frac{s}{t}x\right|^2\right\}
\]
is concave in $x$.  Indeed, the expression in braces is jointly concave in
$(x,y)$; applying that concavity to a pair of maximizing points proves
concavity of its partial supremum.  A finite concave function on $\R^m$ has
an upper supporting affine function at each point.  Adding back the
quadratic therefore gives a vector $q$ such that
\[
  f^\ell(x+h,t)\le f^\ell(x,t)+q\cdot h+\frac{|h|^2}{2t}
  \quad\text{for every }h.
\]
Comparison with \eqref{eq:dp-lower-support}, first for $h=\rho e$ and then
for $h=-\rho e$, and passage to $\rho\downarrow0$, shows that
$q=(y-x)/\ell$.  The two quadratic remainders now prove differentiability
and \eqref{eq:dp-gradient}--\eqref{eq:dp-supports}.
\end{proof}

Fix a state $z=(x,t)$ with $t>\epsilon$.  Choose a maximizer for
$\Penv_{2\epsilon}f(z)$ and a minimizer for $\Menv_{2\epsilon}v(z)$, and
take $x^+$ and $x^-$ to be the respective intermediate points at time
$t+\epsilon$ and $t-\epsilon$ on the straight paths to those optimizers.
The composition formulas for the envelopes and additivity of the quadratic
cost along a straight path give the following identities.  Each half of a
straight path has the same velocity and cost $\epsilon|Q_\pm|^2/2$.
\begin{equation}\label{eq:dp-step-choice}
\begin{aligned}
  f^{2\epsilon}(z)&=f^\epsilon(z^+)-\frac{\epsilon}{2}|Q_+|^2,
       &z^+&=(x^+,t+\epsilon),\\
  \Menv_{2\epsilon}v(z)&=g_\epsilon(z^-)+\frac{\epsilon}{2}|Q_-|^2,
       &z^-&=(x^-,t-\epsilon),
\end{aligned}
\end{equation}
where
\begin{equation}\label{eq:dp-velocities}
  Q_+=\frac{x^+-x}{\epsilon}=\nabla_x f^{2\epsilon}(z),\qquad
  Q_-=\frac{x-x^-}{\epsilon},\qquad Q_0=\nabla_x f^\epsilon(z).
\end{equation}
The formula for $Q_+$ follows from \eqref{eq:dp-gradient} for the full
$2\epsilon$ optimizer.  Although $z^-$ is the backward step, $Q_-$ is
oriented forward from $x^-$ to $x$; the spatial increment for that step is
$-\epsilon Q_-$.  The backward optimizer need not be unique; any choice has
the properties just stated.

\begin{lemma}\label{lem:dp-drift}
At the state and steps above define
\begin{align*}
  D&=f(z)+f^{2\epsilon}(z)-2f^\epsilon(z),\\
  e&=f^\epsilon(z^-)+\frac{\epsilon}{2}|Q_-|^2-f(z),\\
  \widetilde D&=2g_\epsilon(z)-v(z)-\Menv_{2\epsilon}v(z),\\
  \widetilde e&=v(z)+\frac{\epsilon}{2}|Q_+|^2-g_\epsilon(z^+).
\end{align*}
Here $D$ and $\widetilde D$ are the midpoint gaps for forward convexity and
backward concavity.  The quantities $e$ and $\widetilde e$ are the gaps from
using $z$ as a competitor from the opposite step.
All four quantities are nonnegative, and the exact drift identity and the
velocity estimate are
\begin{equation}\label{eq:dp-drift}
\begin{split}
  \frac{J_\epsilon(z^+)+J_\epsilon(z^-)}{2}-J_\epsilon(z)
  &=\frac{D+e+\widetilde D+\widetilde e}{4\epsilon}\\
  &\ge\frac{D+e}{4\epsilon}\\
  &\ge c\bigl(|Q_+-Q_0|^2+|Q_--Q_0|^2\bigr),
  \qquad c=\frac{1}{16(1+\log2)}.
\end{split}
\end{equation}
\end{lemma}

For a fair choice between $z^+$ and $z^-$, the left side of
\eqref{eq:dp-drift} is the conditional expected change of $J_\epsilon$
given this state and the chosen optimizers.
The drift therefore controls the squared discrepancies from the same
reference velocity $Q_0$ at that state.

\begin{proof}
Convexity in the length of $\Penv_\ell f(z)$ gives $D\ge0$, and concavity
in the length of $\Menv_\ell v(z)$ gives $\widetilde D\ge0$.  The point $z$
is a forward competitor from $z^-$ and a backward competitor from $z^+$,
which gives $e\ge0$ and $\widetilde e\ge0$.  The identities in
\eqref{eq:dp-step-choice} and the definitions of the competitor gaps yield
\begin{align*}
  f^\epsilon(z^+)+f^\epsilon(z^-)-2f^\epsilon(z)
    &=D+e+\frac{\epsilon}{2}\bigl(|Q_+|^2-|Q_-|^2\bigr),\\
  g_\epsilon(z^+)+g_\epsilon(z^-)-2g_\epsilon(z)
    &=-\widetilde D-\widetilde e
       +\frac{\epsilon}{2}\bigl(|Q_+|^2-|Q_-|^2\bigr).
\end{align*}
The same difference of the two half-step costs occurs in both lines, so it
cancels on subtraction.  Dividing by $4\epsilon$ proves the identity and
its first inequality.

For the final inequality let $b,d\in\R^m$ be arbitrary.  Apply the
path-convexity property of the forward envelope to
\[
  x(\ell)=x+d+\frac{\ell}{2\epsilon}(b-d),\qquad 0\le\ell\le2\epsilon.
\]
For this path the function
$\Penv_\ell f(x(\ell),t)+|b-d|^2\ell\log\ell/(4\epsilon^2)$ is convex,
with $\ell\log\ell$ assigned its continuous value zero at $\ell=0$.
The midpoint inequality, and
$2\epsilon\log(2\epsilon)-2\epsilon\log\epsilon=2\epsilon\log2$, give
\begin{equation}\label{eq:dp-path-midpoint}
  f(x+d,t)+f^{2\epsilon}(x+b,t)
    -2f^\epsilon\left(x+\frac{b+d}{2},t\right)
  \ge-\frac{\log2}{2\epsilon}|b-d|^2.
\end{equation}
Using $x+d$ as a forward competitor from $z^-$ and expanding its cost gives
\[
  f(x+d,t)\le f(z)+e+Q_-\cdot d+\frac{|d|^2}{2\epsilon}.
\]
Use the upper supporting inequality in \eqref{eq:dp-supports} for
$f^{2\epsilon}(x+b,t)$, and the lower one for
$f^\epsilon(x+(b+d)/2,t)$.  The left side of
\eqref{eq:dp-path-midpoint} is then at most
\[
  D+e+(Q_+-Q_0)\cdot b+(Q_--Q_0)\cdot d
       +\frac{|b|^2}{2t}+\frac{|d|^2}{2\epsilon}
       +\frac{|b+d|^2}{4\epsilon}.
\]
Combining this upper bound with \eqref{eq:dp-path-midpoint}, using
$t>\epsilon$ and $|b\pm d|^2\le2(|b|^2+|d|^2)$, proves
\[
  D+e\ge (Q_0-Q_+)\cdot b+(Q_0-Q_-)\cdot d
            -\frac{1+\log2}{\epsilon}(|b|^2+|d|^2).
\]
Choose
\[
  b=\frac{\epsilon(Q_0-Q_+)}{2(1+\log2)},\qquad
  d=\frac{\epsilon(Q_0-Q_-)}{2(1+\log2)}.
\]
The result is
$D+e\ge\epsilon(|Q_+-Q_0|^2+|Q_--Q_0|^2)/(4(1+\log2))$, as required.
\end{proof}

\subsection{Fitted frames and a logarithmic velocity bound}

The walks will later be viewed in frames based at fits at smaller positive
times.  We first show that the squared discrepancies in
\eqref{eq:dp-drift} are unchanged in those frames, and that the transformed
solution remains in the standing family.  The velocity estimate following
this calculation will then apply with constants independent of the frame.

\begin{lemma}[Covariance in a fitted frame]\label{lem:fitted-frame-covariance}
Let $r>0$, and let $(p_0,a_0)$ be a fit for $v$ at $(x_0,r)$ of radius $r$.
At corresponding points define
\begin{equation}\label{eq:contact-frame}
\begin{gathered}
  s=\frac{t}{r},\qquad y=\frac{x-x_0-p_0(t-r)}{r},\\
  x=x_0+r\bigl(y+p_0(s-1)\bigr),\qquad t=rs,\\
  \bar v(y,s)=\frac{v(x,t)-v(x_0,r)}{r}
          -p_0\cdot y-\left(a_0+\frac{|p_0|^2}{2}\right)(s-1).
\end{gathered}
\end{equation}
Then $\bar v\in\Aclass_{m,C_0}$ is centered at $(0,1)$ and satisfies
\eqref{eq:jensen-bound} with $f=\bar v$, $I=\R$, and the same $\kappa$.  Define
\[
  \bar W(y,s)=\Cenv_s\bar v(y,s)-\frac{\kappa^2s}{2}\qquad(s>0).
\]
For every $\sigma>0$, the barycentric and $W$ gaps satisfy
\begin{equation}\label{eq:contact-W-covariance}
\begin{split}
  \Cenv_\sigma\bar v(y,s)-\bar v(y,s)
       &=\frac{\Cenv_{r\sigma}v(x,t)-v(x,t)}{r},\\
  \bar W(y,s)-\bar v(y,s)&=\frac{W(x,t)-v(x,t)}{r}
       \qquad(s>0).
\end{split}
\end{equation}
For every $\ell>0$, the solution slopes satisfy
\begin{equation}\label{eq:fitted-slope-covariance}
  \Senv_\ell^\pm\bar v(y,s)=\Senv_{r\ell}^\pm v(x,t)-a_0.
\end{equation}
If the steps in \eqref{eq:dp-step-choice} are chosen at any state with
$t>\epsilon$, their images are optimizing steps for the transformed pair,
with step size $\epsilon'=\epsilon/r$.  Their velocities satisfy
\begin{equation}\label{eq:contact-velocity-covariance}
\begin{gathered}
  Q'_+=\frac{y^+-y}{\epsilon'}=Q_+-p_0,\qquad
  Q'_- =\frac{y-y^-}{\epsilon'}=Q_--p_0,\\
  Q'_0=\nabla_y\Penv_{\epsilon'}\bar W(y,s)=Q_0-p_0,\qquad
  Q'_\pm-Q'_0=Q_\pm-Q_0.
\end{gathered}
\end{equation}
All these identities are global at corresponding states, with the indicated
positive-time restriction for $W$; the states need not lie in the fitting
box.  No contact at $(x_0,r)$ is assumed.  The centering and fit assertion
concerns $\bar v$ alone.
\end{lemma}

\begin{proof}
Lemma~\ref{lem:frame-invariance}, followed by the time translation that puts
the fitted point at $(0,1)$, proves the class and centering assertions.  At
each fixed time the spatial offsets are multiplied by $r$, and values are
multiplied by $r$ up to an affine function.  The affine terms cancel from a
barycentric defect, so its quotient by the root mean square radius is
unchanged.  This proves the Jensen assertion.

For the remaining calculations put
\[
  A(y,s)=\frac{v(x_0,r)}{r}+p_0\cdot y
             +\left(a_0+\frac{|p_0|^2}{2}\right)(s-1),
  \qquad \frac{v(x,t)}{r}=\bar v(y,s)+A(y,s).
\]
Scaling the positions of a barycentric configuration and cancelling its
weighted affine terms gives
\[
  \frac{\Cenv_{r\sigma}v(x,t)}{r}
       =\Cenv_\sigma\bar v(y,s)+A(y,s).
\]
Subtracting the displayed formula for $v/r$ proves the first identity in
\eqref{eq:contact-W-covariance}.  For $s>0$, take $\sigma=s$.  The old time
and the old barycentric parameter in $W$ are both $t=rs$, while the
correction $\kappa^2t/(2r)$ is $\kappa^2s/2$.  Hence
\[
  \frac{W(x,t)}{r}=\bar W(y,s)+A(y,s),
\]
which proves the second gap identity.  This is why the frame uses $s=t/r$:
it preserves the time-zero boundary in the definition of $W$.

The quadratic envelopes obey the following value identities for $\ell>0$,
with $s>0$ in the identity involving $W$:
\begin{equation}\label{eq:envelope-covariance}
\begin{split}
  \frac{\Penv_{r\ell}v(x,t)}{r}
    &=\Penv_\ell\bar v(y,s)+A(y,s)+a_0\ell,\\
  \frac{\Menv_{r\ell}v(x,t)}{r}
    &=\Menv_\ell\bar v(y,s)+A(y,s)-a_0\ell,\\
  \frac{\Penv_{r\ell}W(x,t)}{r}
    &=\Penv_\ell\bar W(y,s)+A(y,s)+a_0\ell.
\end{split}
\end{equation}
To check the first identity, parameterize an old forward competitor by a
new position $y'$ at time $s+\ell$.  Its old displacement divided by $r$
is $y'-y+p_0\ell$.  The cost divided by $r$ is
\[
  \frac{|y'-y+p_0\ell|^2}{2\ell}
    =\frac{|y'-y|^2}{2\ell}+p_0\cdot(y'-y)
                                     +\frac{|p_0|^2\ell}{2}.
\]
The affine change in $A$ from $(y,s)$ to $(y',s+\ell)$ is
$p_0\cdot(y'-y)+(a_0+|p_0|^2/2)\ell$, leaving $a_0\ell$ after the cost is
subtracted.  For a backward competitor the old displacement divided by $r$
is $y'-y-p_0\ell$; the same expansion with the cost added leaves
$-a_0\ell$.  These prove the first two identities.  The last follows by the
same forward calculation using $W/r=\bar W+A$ at every positive time.
Subtract the formula for $v/r$ from the forward identity, and subtract the
backward identity from that formula.  Dividing each result by $\ell$ gives
\eqref{eq:fitted-slope-covariance}.

The competitor calculations also show that old optimizers map to optimizers
for the transformed envelopes.  The affine coordinate map sends their
straight intermediate points to the new straight intermediate points.
Thus the mapped steps are precisely the steps of
\eqref{eq:dp-step-choice} for the transformed pair.  The first two velocity
equalities in \eqref{eq:contact-velocity-covariance} follow from
\eqref{eq:contact-frame} and $t^\pm-t=\pm\epsilon$.  For the reference
velocity, differentiate the last identity of \eqref{eq:envelope-covariance}
in $y$ with $r\ell=\epsilon$.  The derivative of the left side is
$\nabla_x\Penv_\epsilon W=Q_0$, since $D_yx=rI$, and the derivative of $A$
is $p_0$.  This proves the equality for $Q'_0$ and the invariance of the
velocity discrepancies, with no additional scale factor.
\end{proof}

It remains to bound the reference velocity in each of these centered
frames.  As before, write $L(K)=1+\log K$ for $K\ge2$.  The next bound is
uniform over the standing family.  Its squared logarithmic singularity in
time is integrable near zero, which will control the expected martingale
variation through the time-walk occupation estimates.

\begin{lemma}\label{lem:dp-log-velocity}
Suppose $K\ge2$, $|x|,|t|\le K$, $0<\epsilon\le1$, and $t>\epsilon$.
Then the velocity in \eqref{eq:dp-velocities} satisfies
\begin{equation}\label{eq:dp-log-velocity}
  |Q_0|=|\nabla_x\Penv_\epsilon W(x,t)|
       \le C\bigl(L(K)+|\log t|\bigr).
\end{equation}
Here $C$ depends only on the fitting constant, the spatial dimension, and
$\kappa$.
\end{lemma}

\begin{proof}
Let $s=t+\epsilon$, so $2\epsilon<s\le K+1$.  At the point $x$, choose any
upper supporting slope $p$ for the semiconcave function $f(\cdot,s)$:
\begin{equation}\label{eq:dp-W-upper-support}
  f(x+h,s)\le f(x,s)+p\cdot h+\frac{|h|^2}{2s}
  \quad(h\in\R^m).
\end{equation}
Such a slope exists because $f(\cdot,s)-|\cdot|^2/(2s)$ is finite and
concave on all of $\R^m$.  If $p\ne0$, insert $h=-sp/|p|$ to obtain
\[
  f(x+h,s)-f(x,s)\le-s|p|+\frac{s}{2}.
\]
The sandwich bound for $W$ gives, in the other direction,
\[
  f(x+h,s)-f(x,s)
    \ge v(x+h,s)-v(x,s)-\frac{\kappa^2s}{2}.
\]
The spatial increment estimate from the fits, applied with increment length
$s$ at the center $(x,s)$ in the box of size $2K$, bounds the right side
below by
\[
  -Cs\bigl(L(2K)+|\log s|\bigr)-\frac{\kappa^2s}{2}.
\]
The resulting estimate for $p$ also holds when $p=0$.  Since
$1<s/t<2$, we have $|\log s|\le|\log t|+\log2$, and hence
\begin{equation}\label{eq:dp-support-slope-bound}
  |p|\le C\bigl(L(K)+|\log t|\bigr).
\end{equation}

Let $h$ be the optimizing displacement for $f^\epsilon(x,t)$.  Its value
is at least the value of the competitor $h=0$.  In conjunction with
\eqref{eq:dp-W-upper-support}, this gives
\[
  \frac{|h|^2}{2\epsilon}\le p\cdot h+\frac{|h|^2}{2s}.
\]
For $h\ne0$, using $s>2\epsilon$ yields $|h|/\epsilon\le4|p|$; the same
inequality is immediate when $h=0$.  By \eqref{eq:dp-gradient},
$Q_0=h/\epsilon$, so \eqref{eq:dp-support-slope-bound} proves the lemma.
\end{proof}
